% ===== PRÉAMBULE CANONIQUE — à coller VERBATIM en tête de CHAQUE .tex =====
\documentclass[11pt,a4paper]{article}
\usepackage[utf8]{inputenc}\usepackage[T1]{fontenc}\usepackage{lmodern}
\usepackage[french]{babel}
\usepackage{amsmath,amssymb,mathtools}\usepackage{icomma}
\usepackage{siunitx}\sisetup{locale=FR}  % \qty{}{}, \si{}, \ang{}
\usepackage{xcolor}\usepackage{colortbl}
\usepackage{tikz}\usepackage{pgfplots}\pgfplotsset{compat=1.18}
\usepackage[siunitx,europeanresistors]{circuitikz} % circuits (élec.)
\usepackage[most]{tcolorbox}\usepackage{enumitem}\usepackage{array,booktabs}
\usepackage[a4paper,margin=2.2cm]{geometry}
\usetikzlibrary{arrows.meta,calc,angles,quotes,positioning,patterns,%
  backgrounds,shapes.geometric,decorations.pathmorphing,decorations.markings,intersections}
\definecolor{primary}{RGB}{222,49,99}\definecolor{primarydark}{RGB}{150,22,55}
\definecolor{defcol}{RGB}{0,114,178}\definecolor{methcol}{RGB}{52,92,148}
\definecolor{excol}{RGB}{0,135,90}\definecolor{attncol}{RGB}{200,40,55}
\tcbset{boxbase/.style={enhanced,breakable,boxrule=0.9pt,arc=2.5pt,
  left=3mm,right=3mm,top=2.3mm,bottom=2.3mm,fonttitle=\bfseries,coltitle=white}}
% --- boîtes TUTORIEL (noms EXACTS, ne pas renommer) ---
\newtcolorbox{cle}[1][]{boxbase,colback=defcol!6,colframe=defcol,
  title={\textsc{À retenir}\ifx\relax#1\relax\relax\else\textmd{ : \itshape #1}\fi}}
\newtcolorbox{methode}[1][]{boxbase,colback=methcol!7,colframe=methcol,sharp corners,
  title={\textsc{Méthode}\ifx\relax#1\relax\relax\else\textmd{ --- \itshape #1}\fi}}
\newtcolorbox{exemplebox}[1][]{boxbase,colback=excol!5,colframe=excol,
  title={\textsc{Exemple}\ifx\relax#1\relax\relax\else\textmd{ : \itshape #1}\fi}}
\newtcolorbox{piege}[1][]{boxbase,colback=attncol!6,colframe=attncol,
  title={\textsc{Piège}\ifx\relax#1\relax\relax\else\textmd{ : \itshape #1}\fi}}
\newtcolorbox{remarque}[1][]{boxbase,colback=black!4,colframe=black!45,
  title={\textsc{Remarque}\ifx\relax#1\relax\relax\else\textmd{ : \itshape #1}\fi}}
% --- boîtes EXERCICE / CORRIGÉ (noms EXACTS, auto-comptées) ---
\newtcolorbox[auto counter]{exo}[1][]{boxbase,colback=primary!4,colframe=primary,
  title={\textsc{Exercice}~\thetcbcounter\ifx\relax#1\relax\relax\else\textmd{ --- \itshape #1}\fi}}
\newtcolorbox[auto counter]{corrige}[1][]{boxbase,colback=excol!4,colframe=excol,
  title={\textsc{Corrigé}~\thetcbcounter\ifx\relax#1\relax\relax\else\textmd{ --- \itshape #1}\fi}}
\newcommand{\vv}[1]{\vec{#1}}\newcommand{\ds}{\displaystyle}
\newcommand{\R}{\mathbb{R}}\newcommand{\N}{\mathbb{N}}\newcommand{\Z}{\mathbb{Z}}
% (un \usepackage{fancyhdr}+en-tête est possible ; il est ignoré côté web)
% =========================================================================
\begin{document}
\begin{corrige}[moment d'une force]
\begin{enumerate}[label=\alph*)]
  \item $\mathcal{M} = F\times b = 30\times 0{,}2 = \qty{6}{N.m}$.
  \item En \textbf{newtons-mètres} (\si{\newton\metre}).
\end{enumerate}
\end{corrige}

\begin{corrige}[plus loin, plus efficace]
\begin{enumerate}[label=\alph*)]
  \item $\mathcal{M}_1 = 30\times 0{,}1 = \qty{3}{N.m}$ ; $\mathcal{M}_2 = 30\times 0{,}3 = \qty{9}{N.m}$.
  \item À la distance la plus grande ($b_2=\qty{0.3}{m}$) : le moment y est le plus élevé, donc le
        boulon tourne le plus facilement. Un grand bras de levier augmente l'effet de rotation.
\end{enumerate}
\end{corrige}

\begin{corrige}[un moment nul]
\begin{enumerate}[label=\alph*)]
  \item La ligne d'action passe par l'axe : la distance perpendiculaire est nulle, donc $b=0$.
  \item $\mathcal{M} = F\times b = F\times 0 = 0$. Le moment est \textbf{nul} : la porte \emph{ne tourne
        pas}. (Pousser vers les gonds ne fait pas pivoter une porte.)
\end{enumerate}
\end{corrige}

\begin{corrige}[moment d'un couple]
\begin{enumerate}[label=\alph*)]
  \item Deux forces égales et opposées : leur somme est \textbf{nulle} ($\vv F+(-\vv F)=\vv 0$). Le
        couple ne provoque aucune translation.
  \item $\mathcal{M} = F\times d = 8\times 0{,}25 = \qty{2}{N.m}$.
\end{enumerate}
\end{corrige}

\begin{corrige}[travail ou moment ?]
\begin{enumerate}[label=\alph*)]
  \item Le \textbf{travail} est une énergie ; il s'exprime en \textbf{joules}
        ($\qty{1}{J}=\qty{1}{N.m}$). Le moment n'est pas une énergie et reste en \si{\newton\metre}.
  \item \textbf{Travail} : la longueur (le déplacement) est \emph{parallèle} à la force.
        \textbf{Moment} : la longueur (le bras de levier) est \emph{perpendiculaire} à la force.
\end{enumerate}
\end{corrige}

\end{document}
